%% ============================================================================
%% sections/10-conclusion.tex -- Conclusion
%% Class-independent section file: it is \input by main-acm.tex and main-ieee.tex.
%% Rules (see README-build.md): no class-specific commands; cite with \cite only;
%% figure sizes come from the float budget and table bodies from tables/, never from this file.
%% Every paragraph and every caption is one numbered unit that starts with a
%% marker line "%% P:<id>".  Ids are frozen: never renumber, never reuse.
%% The same ids appear in paper.zh-CN.md; check_align.py compares the two.
%% Word budget of this section: 100 words of running text.
%% ============================================================================

\section{Conclusion}\label{sec:conclusion}

%% P:sec10-p1
\ptag{sec10-p1}%
This paper asks when a GEMM result should stay \blockscaled across the matrix-to-vector \handoff of a PE, rather than cross as a scalar or be consumed in a fused epilogue, and what a \failvisible \handoff costs.
The placement sweep and the cost breakdown will give the answer in three quantities: whether fusion is feasible, how long the result must live, and the granularity of metadata transactions \TBD{E4}{written after the results}.
That answer is bounded by its evidence: numerical error on synthetic inputs, placements compared within one RTL, and physical costs as \evPostRoute{}s, not silicon measurements.
